\chapter{有限温度 Green 函数与虚时间}\label{chap:thermal-green}

受外力驱动时，量子谐振子在其固有频率处吸收能量。只研究热平衡的振子，也能从位置关联中算出这个频率。\secref{sec:linear-response}用了两个实时时刻的关联。有限温度 Green 函数将两个算符插入同一个热迹，沿虚时间改变它们的间隔，使跃迁能量直接出现在指数衰减率中。

\section{热权重与虚时间演化}
对不显含时间的 Hamilton 量，封闭系统的时间演化算符是$\exp(-\i\hat Ht/\hbar)$，取$t = -\i\tau$便得到$\exp(-\tau\hat H/\hbar)$。在正则系综中，把$\tau$取为
\begin{equation}
	\label{eq:sm-thermal-time-period}
	\tau_T = \beta\hbar = \frac{\hbar}{\kB T}
\end{equation}
就得到热权重$\exp(-\beta\hat H)$。本章的$\tau$具有时间量纲，降温会延长虚时间区间。下文用$\hat K$统一记热权重的生成元，正则系综中$\hat K = \hat H$，满足$[\hat H,\hat N] = 0$的巨正则系综中$\hat K = \hat H - \mu\hat N$。相应地，$\mathcal Z = \tr\e^{-\beta\hat K}$，$\ave{\cdots}_\beta = \tr(\e^{-\beta\hat K}\cdots)/\mathcal Z$。若正则迹限定在固定粒子数的子空间，下文移动算符时要求它保持该子空间。产生、湮灭算符的单粒子关联以下采用对所有粒子数取迹的巨正则系综。

取虚时间算符$\hat A(\tau) = \e^{\tau\hat K/\hbar}\hat A\e^{-\tau\hat K/\hbar}$。对$0 < \tau < \tau_T$，两个算符的热关联为
\begin{equation}
	\ave{\hat A(\tau)\hat B(0)}_\beta
	= \frac1{\mathcal Z}\tr\!\left(
	\e^{-(\tau_T-\tau)\hat K/\hbar}\hat A\,
	\e^{-\tau\hat K/\hbar}\hat B\right).
	\label{eq:sm-thermal-trace}
\end{equation}
若$\hat K\ket{m} = K_m\ket m$，在两个算符之间插入完备集，右边成为
\[
	\sum_{m,n}\frac{\e^{-\beta K_m}}{\mathcal Z}\,
	\e^{-(K_n-K_m)\tau/\hbar}\mel{m}{\hat A}{n}\mel{n}{\hat B}{m}.
\]
初态的热概率决定哪些跃迁贡献显著。对于从基态出发的跃迁，能隙越大，对应项随$\tau$衰减越快。

\section{量子谐振子的位置关联}
考虑质量为$m$、角频率为$\omega_0$的谐振子。用升降算符表示 Hamilton 量和位置，
\begin{gather*}
	\hat H = \hbar\omega_0(\hat a^\dagger\hat a + 1/2),\\
	\hat x = x_0(\hat a + \hat a^\dagger),\qquad
	x_0 = \sqrt{\frac{\hbar}{2m\omega_0}} .
\end{gather*}
对能级求和得到$Z = \e^{-\beta\hbar\omega_0/2}/(1 - \e^{-\beta\hbar\omega_0})$，以及平均激发数$\bar n = \ave{\hat a^\dagger\hat a}_\beta = (\e^{\beta\hbar\omega_0} - 1)^{-1}$。由于$\hat a$与$\hat a^\dagger$各使能量改变一个量子，它们的虚时间演化分别给出因子$\e^{-\omega_0\tau}$和$\e^{\omega_0\tau}$。用$\ave{\hat a\hat a^\dagger}_\beta = \bar n + 1$及$\ave{\hat a\hat a}_\beta = \ave{\hat a^\dagger\hat a^\dagger}_\beta = 0$，得到在$0 < \tau < \tau_T$上的位置关联\footnote{同一振子的虚时间路径积分可参见T.~Senthil的\href{https://ocw.mit.edu/courses/8-513-many-body-theory-for-condensed-matter-systems-fall-2004/resources/notes4/}{MIT 8.513讲义第四篇}。}
\begin{align}
	C_x(\tau) &= \ave{\hat x(\tau)\hat x(0)}_\beta
	= x_0^2\bigb{(\bar n + 1)\e^{-\omega_0\tau} + \bar n\,\e^{\omega_0\tau}}\notag\\
	&= \frac{\hbar}{2m\omega_0}\,
	\frac{\cosh[\omega_0(\tau_T/2 - \tau)]}{\sinh(\omega_0\tau_T/2)} .
	\label{eq:sm-oscillator-imaginary}
\end{align}
第一项由$\hat x(0)$将振子升高一个能级得到。第二项要求初态里已有激发，降到一个较低能级后便给出符号相反的指数。两项相加使$C_x(\tau_T - \tau) = C_x(\tau)$。

取$\tau \to 0^+$得到位置的均方涨落
\begin{equation}
	\label{eq:sm-oscillator-thermal-variance}
	\ave{\hat x^2}_\beta = \frac{\hbar}{2m\omega_0}
	\coth\!\Bigp{\frac{\beta\hbar\omega_0}{2}} .
\end{equation}
在$\kB T \gg \hbar\omega_0$时还原出能量均分结果$\ave{\hat x^2} = \kB T/(m\omega_0^2)$。在$T \to 0$时，等时涨落仍有基态的宽度$x_0^2$，固定$\tau > 0$的关联趋于$x_0^2\e^{-\omega_0\tau}$，其衰减率量出了从基态到第一激发态的能隙$\hbar\omega_0$。

\section{热迹的循环性与Matsubara频率}
振子公式在虚时间区间的两端取相同的值。这是热迹首尾相接的一个例子。对一般的式\eqref{eq:sm-thermal-trace}，用迹的循环性移动算符并合并指数因子，得到
\begin{equation}
	\label{eq:sm-kms-relation}
	\ave{\hat A(\tau)\hat B(0)}_\beta
	= \ave{\hat B(0)\hat A(\tau - \tau_T)}_\beta .
\end{equation}
这称为热平衡的Kubo--Martin--Schwinger条件，简称KMS条件。迹的循环性对 Bose 和 Fermi 算符都成立。为了把正负$\tau$写成一个函数，用虚时间排序$\mathcal T_\tau$把较晚的算符移到左侧。交换两个 Fermi 奇算符时取$\eta = -1$，对于 Bose 算符取$\eta = +1$。例如，Fermi 子的数密度算符$\hat\psi^\dagger\hat\psi$是两个 Fermi 算符的乘积，遵守 Bose 型排序。定义
\begin{equation}
	\label{eq:sm-imaginary-time-correlation}
	C_{AB}(\tau) = \ave{\mathcal T_\tau\hat A(\tau)\hat B(0)}_\beta
	= \begin{cases}
		\ave{\hat A(\tau)\hat B(0)}_\beta,&0 < \tau < \tau_T,\\
		\eta\ave{\hat B(0)\hat A(\tau)}_\beta,&-\tau_T < \tau < 0 .
	\end{cases}
\end{equation}
由KMS条件可将其延拓为$C_{AB}(\tau + \tau_T) = \eta C_{AB}(\tau)$。Fermi 关联的反周期负号来自排序中的一次交换。Fermi 子的单粒子 Green 函数通常还在$C_{AB}$的定义前加一个总负号，这不影响边界条件。

在一个热周期内把有序关联展开为 Fourier 级数，
\begin{subequations}\label{eq:sm-matsubara-transform}
	\begin{align}
		C_{AB}(\tau) &= \frac1{\tau_T}\sum_n C_{AB}(\i\nu_n)\e^{-\i\nu_n\tau},\\
		C_{AB}(\i\nu_n) &= \int_0^{\tau_T}C_{AB}(\tau)\e^{\i\nu_n\tau}\d\tau .
	\end{align}
\end{subequations}
这里$\i\nu_n$表示复频率轴上的点。由$\e^{-\i\nu_n\tau_T} = \eta$，允许的角频率（Matsubara频率）为
\begin{equation}
	\label{eq:sm-matsubara-frequencies}
	\nu_n = \begin{cases}
		2\pi n/(\beta\hbar),&\text{Bose 型},\\
		(2n + 1)\pi/(\beta\hbar),&\text{Fermi 型},
	\end{cases}
	\qquad n \in \mathbb Z .
\end{equation}
Bose 频率允许零模$\nu_0 = 0$，Fermi 频率的最小绝对值为$\pi/(\beta\hbar)$。

以能量为$\varepsilon$的一个 Fermi 轨道为例，令$\xi = \varepsilon - \mu$，$n_\Fermi = (\e^{\beta\xi} + 1)^{-1}$。湮灭算符的演化为$\hat c(\tau) = \e^{-\xi\tau/\hbar}\hat c$。由于$\ave{\hat c\hat c^\dagger} = 1 - n_\Fermi$、$\ave{\hat c^\dagger\hat c} = n_\Fermi$，该轨道的单粒子 Green 函数为
\begin{equation}
	\label{eq:sm-fermion-imaginary-time-green}
	G(\tau) = -\ave{\mathcal T_\tau\hat c(\tau)\hat c^\dagger(0)}_\beta
	= \begin{cases}
		-(1 - n_\Fermi)\e^{-\xi\tau/\hbar},&0 < \tau < \tau_T,\\
		n_\Fermi\,\e^{-\xi\tau/\hbar},&-\tau_T < \tau < 0 .
	\end{cases}
\end{equation}
$G(0^-) = n_\Fermi$直接给出占据数。两侧的跳变$G(0^+) - G(0^-) = -1$来自 Fermi 反对易关系，而$(1 - n_\Fermi)\e^{-\beta\xi} = n_\Fermi$保证了反周期性。

\section{从虚频关联到实频响应}
将振子的式\eqref{eq:sm-oscillator-imaginary}按 Bose 频率变换。用$\e^{\i\nu_n\tau_T} = 1$分别积分两个指数项，可得
\begin{align}
	C_x(\i\nu_n)
	&= x_0^2\biggp{\frac1{\omega_0 - \i\nu_n} + \frac1{\omega_0 + \i\nu_n}}
	= \frac{\hbar}{m(\omega_0^2 + \nu_n^2)} .
\end{align}
现在给振子施加力$f(t)$，耦合能为$-f(t)\hat x$。未受扰动的位置算符随实时间的演化为
\[
	\hat x(t) = \hat x(0)\cos\omega_0t + \frac{\hat p(0)}{m\omega_0}\sin\omega_0t .
\]
由$[\hat x,\hat p] = \i\hbar$以及Kubo公式\eqref{eq:sm-kubo-retarded}，
\begin{subequations}
	\label{eq:sm-oscillator-retarded-response}
	\begin{align}
		\chi^R_{xx}(t) &= \Theta(t)\frac{\sin\omega_0t}{m\omega_0},\\
		\chi^R_{xx}(\omega)
		&= \lim_{\epsilon\to0^+}\frac1{m[\omega_0^2 - (\omega + \i\epsilon)^2]} .
	\end{align}
\end{subequations}
因此$C_x(\i\nu_n)/\hbar$在正虚频处和$\chi^R_{xx}(\omega)$在实轴上缘，都是同一个解析函数$1/[m(\omega_0^2 - z^2)]$在不同位置的取值。对这个振子，虚频零点给出静力下的位移与力之比$1/(m\omega_0^2)$，实轴的极点给出驱动振子的共振频率。

对其他体系也可从能量本征态的谱表示作解析延拓。巨正则系综的虚时间由$\hat K = \hat H - \mu\hat N$生成，Fermi 子的单粒子 Green 函数的极点因此给出相对于化学势的加粒子或去粒子能量。对保持粒子数的观测量，跃迁两端的$N$相同，能量差中的化学势项抵消。

零模还可能含有静态涨落。例如，守恒粒子数的偏差$\delta\hat N = \hat N - \ave{\hat N}$在巨正则系综中具有与$\tau$无关的关联$\ave{(\delta\hat N)^2}$，完全落在$\nu_n = 0$，而$[\hat N(t),\hat N(0)] = 0$，相应的推迟响应为零。数值计算若只给出有限个带误差的虚频数据，还需额外的谱信息或精度约束才能可靠地分辨实频峰。
